% Results. The brief: qualitative results that show the behaviour and its limits, plus quantitative
% results against conventional or alternative approaches. Marking needs at least one plot with axes,
% units and an uncertainty estimate. Follow the order of Table 1 and quote the metric and baseline for each.
% Every figure and table here must be referred to and explained in the text, and every caption must
% make sense on its own.
%
% Planned figures:
%  - Panel against card: PSFs of the open disc, Zhou A and Zhou B side by side, with the NCC under each.
%  - Pair spectra: shared-zero fraction for the ideal, card and panel pairs and the Lab 3 pair, with error bars.
%  - Depth: estimated against true distance for the pair and each single pattern, error bars over repeats.
%    This is the candidate key figure for the repository script.
%  - Refocus: edge width against refocus depth for the 3 x 3 and 5 x 5 scans, with the b/N prediction.
%  - Bokeh: NCC against field angle for the panel star and the card star.
%  - Flutter: MSE ratio against light level for both baselines, plus one deblurred example.
\section{Results}\label{sec:results}

\subsection{Panel characterisation}\label{sec:res-panel}
\tbd{To write.}

\subsection{Depth from a coded pair}\label{sec:res-depth}
\tbd{To write.}

\subsection{Post-capture refocus}\label{sec:res-refocus}
\tbd{To write.}

\subsection{Custom bokeh}\label{sec:res-bokeh}
\tbd{To write.}

\subsection{Flutter shutter}\label{sec:res-flutter}
\tbd{To write.}
